\documentclass[a4paper,12pt]{article}
\usepackage[T2A]{fontenc}
\usepackage[utf8]{inputenc}
\usepackage[russian]{babel}
\usepackage{amsmath,amssymb,mathtools}
\usepackage{icomma}
\usepackage[a4paper,left=24mm,right=24mm,top=24mm,bottom=25mm,headheight=25pt]{geometry}
\usepackage{microtype}
\usepackage{tikz}
\usetikzlibrary{arrows.meta,positioning,shapes.geometric,calc}
\usepackage{pgfplots}
\pgfplotsset{compat=1.18}
\usepackage{enumitem}
\usepackage{tabularx,booktabs}
\usepackage{xcolor}
\usepackage{fancyhdr}
\usepackage{setspace}
\usepackage{hyperref}
\usepackage{needspace}
\usepackage{parskip}
\onehalfspacing
\definecolor{accent}{HTML}{008080}
\definecolor{darkaccent}{HTML}{005B5B}
\definecolor{textgray}{HTML}{353B3E}
\definecolor{muted}{HTML}{596368}
\hypersetup{colorlinks=true,linkcolor=darkaccent,urlcolor=darkaccent,pdfauthor={Лёвин Артём Александрович},pdftitle={Чек-лист: логарифмические уравнения и графики}}
\setlength{\parindent}{0pt}
\setlist[enumerate]{leftmargin=*,itemsep=0.55em,topsep=0.3em}
\pagestyle{fancy}
\fancyhf{}
\fancyhead[L]{\small\color{muted}Логарифмические уравнения и графики}
\fancyhead[R]{\small\color{muted}Ярослав Гаврилов}
\fancyfoot[C]{\small\color{muted}\thepage}
\renewcommand{\headrulewidth}{0.4pt}
\renewcommand{\headrule}{\hbox to\headwidth{\color{accent}\leaders\hrule height \headrulewidth\hfill}}
\newcommand{\topic}[1]{\vspace{0.75em}{\large\color{darkaccent}#1}\par\vspace{0.3em}{\color{accent}\rule{\linewidth}{0.35pt}}\vspace{0.4em}}
\newcommand{\subtopic}[1]{\par\vspace{0.45em}{\color{darkaccent}\textbf{#1}}\par\vspace{0.15em}}
\newcommand{\key}[1]{{\color{darkaccent}\textbf{#1}}}
\newcommand{\warning}[1]{{\color{darkaccent}\textbf{Обратите внимание.}} #1}
\begin{document}
\thispagestyle{empty}
\begin{center}
\vspace*{28mm}
{\large\color{darkaccent}МАТЕМАТИКА}\par
\vspace{25mm}
{\LARGE\color{textgray}Чек-лист}\par
\vspace{10mm}
{\Large\color{accent}Логарифмические уравнения,\\[3pt] прикладные задачи и графики}\par
\vspace{13mm}
{\normalsize\color{textgray}Подготовка к ЕГЭ по профильной математике}\par
\vfill
{\normalsize Ярослав Гаврилов}\par
\vspace{10mm}
{\normalsize 29 октября 2026 г.}\par
\vspace{21mm}
{\small Лёвин Артём Александрович}\par
{\small эксклюзивно для Ярослава Гаврилова}\par
\vspace{18mm}
\begin{tikzpicture}[x=1cm,y=1cm]
\draw[accent,line width=0.8pt] (0,0) .. controls (1.0,0.55) and (2.2,-0.45) .. (3.5,0)
  .. controls (4.6,0.55) and (5.9,-0.45) .. (7.1,0);
\end{tikzpicture}
\end{center}
\clearpage
\setcounter{page}{1}

\topic{1. Определение и область допустимых значений}
\textbf{Логарифм} числа $b$ по основанию $a$ --- показатель степени, в которую нужно возвести $a$, чтобы получить $b$:
\[
\boxed{\log_a b=c\quad\Longleftrightarrow\quad a^c=b.}
\]
\key{Приём «циклоп»:} из записи $\log_a b=c$ получите $b=a^c$; основание $a$ возводится в степень $c$ и даёт аргумент $b$.

\subtopic{Три обязательных условия}
\[
\boxed{a>0,\qquad a\ne1,\qquad b>0.}
\]
Если аргумент содержит переменную, его положительность нужно проверять для всех найденных корней. Эти условия называют \textbf{областью допустимых значений} (ОДЗ).

\subtopic{Пример 1. Простейшее уравнение}
\[
\log_3(x+1)=2\ \Longleftrightarrow\ x+1=3^2=9\ \Longleftrightarrow\ x=8.
\]
ОДЗ: $x+1>0$, следовательно, $8$ подходит. \key{Ответ: $8$.}

\topic{2. Когда можно убрать логарифмы}
Если основания одинаковы и обе части состоят из одного логарифма, то
\[
\boxed{\log_a U=\log_a V\ \Longleftrightarrow\ U=V,\quad U>0,\ V>0.}
\]
Если перед логарифмом есть коэффициент, рядом находится дополнительное слагаемое или основания различны, сначала преобразуйте выражения.

\subtopic{Пример 2. Равенство логарифмов}
\[
\log_2(x+1)=\log_2(2x+1).
\]
ОДЗ: $x+1>0$, $2x+1>0$. После сравнения аргументов:
\[
 x+1=2x+1\ \Longleftrightarrow\ x=0.
\]
Проверка: $0+1>0$ и $2\cdot0+1>0$. \key{Ответ: $0$.}

\topic{3. Основные свойства логарифмов}
Для $U>0$, $V>0$, $a>0$, $a\ne1$:
\[
\begin{aligned}
\log_a(UV)&=\log_aU+\log_aV,\\
\log_a\frac{U}{V}&=\log_aU-\log_aV,\\
\log_a(U^r)&=r\log_aU,\\
\log_a1&=0,\qquad \log_aa=1.
\end{aligned}
\]
Последнее свойство позволяет заменить число логарифмом: например, $1=\log_2 2$.

\vspace{1.5em}
\topic{4. Дополнительное слагаемое и коэффициент}
\subtopic{Пример 3. Число справа от логарифма}
\[
\log_2(x+1)=\log_2(2x+1)+1.
\]
\textbf{Шаг 1.} Запишите ОДЗ: $x+1>0$, $2x+1>0$.

\textbf{Шаг 2.} Представьте $1$ как $\log_2 2$, затем объедините логарифмы:
\[
\begin{aligned}
\log_2(x+1)&=\log_2(2x+1)+\log_2 2,\\
\log_2(x+1)&=\log_2\bigl(2(2x+1)\bigr).
\end{aligned}
\]
\textbf{Шаг 3.} Теперь основания и форма обеих частей совпадают:
\[
 x+1=4x+2\ \Longleftrightarrow\ -3x=1\ \Longleftrightarrow\ x=-\frac13.
\]
Проверка: $x+1=\frac23>0$ и $2x+1=\frac13>0$. \key{Ответ: $-\frac13$.}

\Needspace{10\baselineskip}
\subtopic{Пример 4. Множитель перед логарифмом}
\[
\log_5(5-x)=2\log_5 3.
\]
Используйте $2\log_5 3=\log_5(3^2)=\log_5 9$:
\[
5-x=9\ \Longleftrightarrow\ x=-4.
\]
ОДЗ: $5-x>0$. При $x=-4$ аргумент равен $9>0$. \key{Ответ: $-4$.}

\topic{5. Основание в виде степени}
При допустимых аргументах и $k\ne0$:
\[
\boxed{\log_{a^k}b=\frac1k\log_ab,\qquad a>0,\ a\ne1,\ b>0.}
\]
\subtopic{Пример 5. Основание $8=2^3$}
\[
\log_8\bigl(2^{8x-4}\bigr)=4.
\]
Показательное выражение $2^{8x-4}$ положительно для всех действительных $x$. Поэтому ОДЗ выполняется автоматически:
\[
\begin{aligned}
\log_{2^3}\bigl(2^{8x-4}\bigr)&=\frac13\log_2\bigl(2^{8x-4}\bigr),\\
\frac{8x-4}{3}&=4\ \Longleftrightarrow\ 8x-4=12\ \Longleftrightarrow\ x=2.
\end{aligned}
\]
\key{Ответ: $2$.} Здесь важно сохранить скобки вокруг показателя $8x-4$.

\Needspace{12\baselineskip}
\vspace{1.5em}
\topic{6. Разные основания: приводим к одному}
\subtopic{Пример 6. Основание $0{,}5$}
\[
\log_2\frac{x}{6}=\log_{0{,}5}(x+1).
\]
ОДЗ: $\frac{x}{6}>0$, $x+1>0$; вместе получаем $x>0$.

Так как $0{,}5=2^{-1}$, то $\log_{0{,}5}(x+1)=-\log_2(x+1)$. Поэтому
\[
\begin{aligned}
\log_2\frac{x}{6}&=-\log_2(x+1),\\
\log_2\frac{x}{6}&=\log_2\frac1{x+1},\\
\frac{x}{6}&=\frac1{x+1},\\
x(x+1)&=6,\\
x^2+x-6&=0,\\
(x+3)(x-2)&=0.
\end{aligned}
\]
Получены кандидаты $x=-3$ и $x=2$. По ОДЗ $x>0$, поэтому $-3$ исключаем. \key{Ответ: $2$.}

\topic{7. Что делать, если встретилось выражение со степенью}
Для любого действительного $t$ число $2^t>0$. Поэтому аргумент $2^{8x-4}$ допустим при любом $x$. Основание степени должно быть положительным, если рассматривается произвольный действительный показатель.

\subtopic{Частые ошибки}
\begin{enumerate}[itemsep=0.12em,topsep=0.12em]
\item Убирать логарифмы, когда остаётся слагаемое вне логарифма.
\item Применять $\log_a U=\log_b V\Rightarrow U=V$ при $a\ne b$.
\item Забывать условия $U>0$ и $V>0$ и записывать посторонний корень.
\item Записывать $\log_a(U+V)=\log_aU+\log_aV$: такой формулы нет.
\item Терять скобки в $r\log_a U$ и при переносе степени: $\log_2(2^{8x-4})=8x-4$.
\item Путать основание с показателем: $\log_{8}\!\left(2^{8x-4}\right)$.
\end{enumerate}

\Needspace{14\baselineskip}
\vspace{1.5em}
\topic{8. Прикладная задача с логарифмом}
В задачах этого типа формула дана в условии. Ваша задача --- аккуратно выписать данные, подставить их, упростить числовой коэффициент и перейти от логарифма к степени.

Рассмотрим формулу
\[
 Q=\frac{\alpha c m}{\gamma}\log_2\frac{T_1-T_0}{T-T_0}.
\]
По данным $\alpha=0{,}7$, $c=4200$, $m=0{,}4$, $\gamma=420$, $Q=8{,}4$, $T_1=100$, $T_0=20$ требуется найти $T$.

\subtopic{Шаг 1. Коэффициент удобнее вычислить отдельно}
\[
\begin{aligned}
\frac{\alpha cm}{\gamma}
&=\frac{0{,}7\cdot4200\cdot0{,}4}{420}
=\frac{7\cdot42\cdot4}{42}\cdot10^{-1}=2{,}8.
\end{aligned}
\]
Здесь можно сократить одинаковые числовые множители до умножения. Равносильный путь --- перевести десятичные дроби в степени десяти.

\subtopic{Шаг 2. Подставьте данные и выделите логарифм}
\[
8{,}4=2{,}8\log_2\frac{80}{T-20}
\quad\Longleftrightarrow\quad
\log_2\frac{80}{T-20}=3.
\]
\subtopic{Шаг 3. Перейдите к степени}
\[
\frac{80}{T-20}=2^3=8
\quad\Longleftrightarrow\quad
T-20=10
\quad\Longleftrightarrow\quad T=30.
\]
\subtopic{Шаг 4. Проверка ОДЗ}
Аргумент $\frac{80}{T-20}$ должен быть положительным. При $T=30$ он равен $8>0$. Исходная формула даёт $2{,}8\cdot3=8{,}4$. \key{Ответ: $T=30$.}

\warning{В прикладных формулах сначала выясняйте смысл величин и их единицы измерения. Сокращать можно лишь математически одинаковые множители; результат должен иметь допустимый физический смысл.}

\Needspace{13\baselineskip}
\vspace{1.5em}
\topic{9. Логарифмическая функция и метод узловых точек}
Рассмотрим функцию
\[
\boxed{f(x)=\log_a x+b,\qquad a>0,\ a\ne1,\ x>0.}
\]
Неизвестных коэффициентов два: $a$ и $b$. Поэтому по \textbf{двум точкам графика} составляем систему из двух уравнений. Особенно удобна точка с абсциссой $x=1$: $\log_a1=0$, откуда сразу $f(1)=b$.

\subtopic{Пример 7. Восстановление формулы по двум точкам}
График проходит через $A(1;-3)$ и $B(4;-1)$. Требуется найти $f(32)$.
\[
\begin{cases}
-3=\log_a1+b,\\
-1=\log_a4+b.
\end{cases}
\quad\Longrightarrow\quad
b=-3,\quad \log_a4=2.
\]
Перейдём к степени: $a^2=4$. Формально $a=\pm2$, но основание логарифма обязано быть положительным, поэтому $a=2$.
\[
 f(x)=\log_2x-3,\qquad f(32)=\log_2 32-3=5-3=\boxed{2}.
\]
\begin{center}
\begin{tikzpicture}
\begin{axis}[
 width=0.78\linewidth,height=64mm,
 axis lines=middle,xmin=0,xmax=8.6,ymin=-6.2,ymax=0.6,
 xtick={1,2,4,8},ytick={-5,-4,-3,-2,-1,0},
 xlabel={$x$},ylabel={$y$},xlabel style={at={(axis description cs:1,0.05)},anchor=east},
 ylabel style={at={(axis description cs:0.055,1)},anchor=north},
 tick label style={font=\scriptsize},label style={font=\small},
 axis line style={darkaccent},grid=major,grid style={gray!15},
 clip=false]
\addplot[accent,thick,domain=0.08:8.5,samples=180]{ln(x)/ln(2)-3};
\addplot[only marks,mark=*,mark size=2.1pt,darkaccent] coordinates {(1,-3) (4,-1)};
\node[above right,font=\footnotesize,text=darkaccent] at (axis cs:1,-3) {$A(1;-3)$};
\node[below right,font=\footnotesize,text=darkaccent] at (axis cs:4,-1) {$B(4;-1)$};
\end{axis}
\end{tikzpicture}
\end{center}
\vspace{-0.7em}
\warning{Точка с $x=0$ на график логарифмической функции не входит. Если получаете $a=-2$, исключите это значение по ограничению $a>0$.}

\clearpage
\topic{10. Алгоритм решения логарифмического уравнения}
\vspace{5mm}
\begin{center}
\begin{tikzpicture}[
  actionbox/.style={draw=accent!85!black,rounded corners=2pt,text width=11.1cm,align=center,minimum height=1.6cm,font=\normalsize,inner sep=6pt},
  terminal/.style={draw=accent!85!black,ellipse,text width=6cm,align=center,minimum height=1cm,inner sep=4pt},
  arrow/.style={-Latex,line width=0.85pt,color=darkaccent}
]
\node[terminal] (start) {Дано логарифмическое уравнение};
\node[actionbox,below=10mm of start] (odz) {Запишите ОДЗ: основания положительны и не равны $1$; аргументы логарифмов строго положительны};
\node[actionbox,below=10mm of odz] (transform) {Упростите запись: перенесите коэффициенты в показатели, числа представьте логарифмами, приведите основания к одному виду};
\node[actionbox,below=10mm of transform] (combine) {Если получили $\log_a U=\log_a V$, приравняйте аргументы; если получили $\log_a U=c$, используйте $U=a^c$};
\node[actionbox,below=10mm of combine] (solve) {Решите полученное уравнение школьным способом};
\node[actionbox,below=10mm of solve] (check) {Проверьте каждый найденный корень по ОДЗ и подстановкой, если это требуется};
\node[terminal,below=10mm of check] (finish) {Запишите допустимые корни};
\draw[arrow] (start)--(odz);
\draw[arrow] (odz)--(transform);
\draw[arrow] (transform)--(combine);
\draw[arrow] (combine)--(solve);
\draw[arrow] (solve)--(check);
\draw[arrow] (check)--(finish);
\end{tikzpicture}
\end{center}
\clearpage

\topic{Тренировочный блок \;\textnormal{(Путь А)}}
Во всех уравнениях указывайте ОДЗ и проверяйте найденные корни. Задания расположены по возрастанию сложности.

\subtopic{Базовый уровень}
\begin{enumerate}[label=\textbf{\arabic*.},itemsep=0.5em]
\item Решите уравнение: $\log_3(x-2)=2$.
\item Решите уравнение: $\log_2(x+5)=\log_2(3x-1)$.
\item Решите уравнение: $\log_3(x+4)=\log_3(2x+1)+1$.
\end{enumerate}

\subtopic{Уровень выше среднего}
\begin{enumerate}[label=\textbf{\arabic*.},start=4,itemsep=0.5em]
\item Решите уравнение: $\log_5(12-x)=2\log_5 2$.
\item Решите уравнение: $\log_9\bigl(3^{4x-2}\bigr)=3$.
\item Решите уравнение: $\log_2\dfrac{x}{3}=\log_{0{,}5}(x+2)$.
\item В модели $Q=k\log_2\dfrac{T_1-T_0}{T-T_0}$ известны $Q=6$, $k=2$, $T_1=100$, $T_0=20$. Найдите $T$.
\end{enumerate}

\subtopic{Повышенный уровень}
\begin{enumerate}[label=\textbf{\arabic*.},start=8,itemsep=0.55em]
\item По формуле $Q=\dfrac{\alpha cm}{\gamma}\log_2\dfrac{T_1-T_0}{T-T_0}$ найдите $T$, если $Q=8$, $\alpha=0{,}6$, $c=3600$, $m=0{,}5$, $\gamma=540$, $T_1=116$, $T_0=20$.
\item График функции $f(x)=\log_a x+b$ проходит через точки $(1;-2)$ и $(9;0)$. Найдите $f(27)$.
\item График функции $f(x)=\log_a x+b$ проходит через точки $(1;4)$ и $\bigl(\frac18;1\bigr)$. Найдите $f(32)$.
\end{enumerate}

\clearpage
\topic{Ответы к тренировочному блоку}
\begin{center}
\renewcommand{\arraystretch}{1.45}
\begin{tabular}{@{}c>{\raggedright\arraybackslash}p{9.7cm}@{}}
\toprule
\textbf{№} & \textbf{Ответ} \\
\midrule
1 & $x=11$ \\
2 & $x=3$ \\
3 & $x=\dfrac15$ \\
4 & $x=8$ \\
5 & $x=2$ \\
6 & $x=1$ \\
7 & $T=30$ \\
8 & $T=26$ \\
9 & $f(27)=1$ \\
10 & $f(32)=9$ \\
\bottomrule
\end{tabular}
\end{center}

\vspace{8mm}
\topic{Итоговая памятка}
\renewcommand{\arraystretch}{1.55}
\begin{tabularx}{\linewidth}{@{}p{0.40\linewidth}X@{}}
\toprule
\textbf{Сигнал в задаче} & \textbf{Полезное действие} \\
\midrule
$\log_a U=c$ & Перейдите к $U=a^c$. \\
$\log_a U=\log_a V$ & Приравняйте аргументы и проверьте ОДЗ. \\
Слагаемое вне логарифма & Представьте его логарифмом или перенесите. \\
Коэффициент перед логарифмом & Используйте $r\log_a U=\log_a(U^r)$ при $U>0$. \\
Разные основания & Приведите их к общей основе. \\
Прикладная формула & Выпишите данные, упростите коэффициент, изолируйте логарифм. \\
$f(x)=\log_a x+b$ & Подставьте координаты двух точек; при $x=1$ сразу $b=f(1)$. \\
\bottomrule
\end{tabularx}
\end{document}
